% Native semantic reconstruction of source physical page 12 (printed page 11).
% Provenance: IEEE Std 1516-2025, document opening and Clause 1.1--1.3.
\HLAStartSemanticOpenerPage{12}{11}
\HLAContentMarker{clause-1}
\fontsize{10.02pt}{11.523pt}\selectfont
\begingroup
\hbadness=10000\relax
\addfontfeatures{FakeStretch=1.035}
\vspace*{65.51pt}

{\HLAFacSans\bfseries\fontsize{19.98pt}{23.02pt}\selectfont
IEEE Standard for Modeling and Simulation\linebreak
(M\&S) High Level Architecture (HLA)---\linebreak
Framework and Rules\par}
\vspace{63.80pt}

\HLAClauseHeading{clause-1}{1. Overview}
\vspace{24.06pt}

\HLASubclauseHeading{clause-1.1}{1.1 Scope}
\vspace{12.62pt}

This document provides an overview of the High Level Architecture (HLA), defines a family of related\linebreak
HLA documents, and defines the principles of HLA in terms of responsibilities that federates (simulations,\linebreak
supporting utilities, or interfaces to live systems) and federations (sets of federates working together)\linebreak
upholds.\par
\vspace{21.06pt}

\HLASubclauseHeading{clause-1.2}{1.2 Purpose}
\vspace{12.62pt}

Many different classes of simulations exist. Each class has changing application characteristics and needs\linebreak
to be flexibly supported to allow for interoperability and reuse across classes and to limit the need to\linebreak
maintain multiple interoperability approaches. The HLA is a common integrated architecture that has been\linebreak
developed to provide a flexible approach to addressing these interoperability and reuse needs. This\linebreak
document describes the general principles defining the HLA and delineates the set of rules that apply to\linebreak
HLA federations and federates. Each rule is then described, and the rationale for its inclusion is provided.\par
\vspace{21.09pt}

\HLASubclauseHeading{clause-1.3}{1.3 Word usage}
\vspace{12.62pt}

The word \textit{shall} indicates mandatory requirements strictly to be followed in order to conform to the standard\linebreak
and from which no deviation is permitted (\textit{shall} equals \textit{is required to}).\textsuperscript{6,7}\par
\vspace{12.10pt}

The word \textit{should} indicates that among several possibilities one is recommended as particularly suitable,\linebreak
without mentioning or excluding others\HLASemicolon{} or that a certain course of action is preferred but not necessarily\linebreak
required (\textit{should} equals \textit{is recommended that}).\par
\vspace{10.06pt}

The word \textit{may} is used to indicate a course of action permissible within the limits of the standard (\textit{may}\linebreak
equals \textit{is permitted to}).\par
\vspace{10.08pt}

The word \textit{can} is used for statements of possibility and capability, whether material, physical, or causal (\textit{can}\linebreak
equals \textit{is able to}).\par

\footnotetext[6]{The use of the word \textit{must} is deprecated and cannot be used when stating mandatory requirements, \textit{must} is used only to describe unavoidable situations.}
\footnotetext[7]{The use of \textit{will} is deprecated and cannot be used when stating mandatory requirements, \textit{will} is only used in statements of fact.}
\endgroup
\clearpage
